\chapter{绪论}

\section{研究背景与意义}
% 交通资产清查的现状（人工成本高、小目标易遗漏、损毁资产难以识别）、智慧交通的发展需求。
\subsection{研究背景}
随着全球交通基础设施规模的持续扩张，如何实现高效、精细和可持续的资产管理，已成为建设“交通强国”和推进交通数字化转型的重要议题。《中华人民共和国国民经济和社会发展第十四个五年规划和2035年远景目标纲要》明确提出要加快交通等传统基础设施的数字化改造，在智能交通等领域开展试点示范 \cite{china2021plan}；交通运输部发布的《公路数字化专项行动方案（2023—2027年）》进一步强调，要围绕公路全生命期管理构建“一套模型、一套数据”的闭环体系，为交通资产清查、数字台账更新和养护决策提供数据基础 \cite{mot2023highway}。在这一背景下，交通资产清查已不再是简单的现场盘点工作，而是支撑交通基础设施精细化治理和智能运维的关键基础环节。

然而，现阶段交通资产清查与管养实践与上述目标之间仍存在明显差距。长期以来，交通资产巡检高度依赖人工观察与人工记录，这种模式不仅人力成本高、效率低，而且容易受到视觉疲劳、主观经验差异与作业环境复杂性的影响。相关研究表明，人工巡检在常规工业视觉检测任务中的平均漏检率约为 $15\%$，在高压或极度疲劳状态下甚至可能升至 $40\%$ \cite{naqviquality}。对于高速公路、城市快速路和复杂道路节点而言，人工方式还面临作业风险高、连续性差和数据标准化程度不足等问题，难以满足大规模、常态化交通资产清查的现实需求 \cite{roadvision_ai}。

从技术角度看，面向交通资产清查的智能感知仍受两类核心瓶颈制约。其一是\textbf{资产发现不完整}：道路场景中的路灯、标志牌、井盖、锥桶及其局部损伤往往具有小目标、远距离、边缘分布和密集分布等特点，传统单次全图检测器在复杂背景和尺度压缩条件下容易发生漏检 \cite{small_object2024}。其二是\textbf{异常状态记录不充分}：交通资产的锈蚀、破损、弯曲、移位和满溢等异常状态具有明显长尾特征，真实采集样本稀缺且分布极不均衡，导致现有识别模型在低频异常属性上的泛化能力不足 \cite{longtail_driving2024}。前者决定系统能否“找到资产”，后者决定系统能否“看懂状态”，二者共同影响交通资产清查结果的完整性、准确性和可用性。

随着视觉语言模型（Vision-Language Models, VLM）、多智能体系统和生成式 AI 的快速发展，交通感知正由单次被动识别逐步走向任务驱动的复合感知新范式。一方面，大模型具备一定的推理、调度和多模态理解能力，可用于组织“先整体观察、再重点复查”的主动感知流程；另一方面，行业标准知识、专家经验与生成模型相结合，也为长尾异常属性样本的可控补充提供了新路径 \cite{highway_scene2025}。因此，围绕交通资产清查任务构建“主动感知检测 + 知识引导属性识别增强”的方法体系，具有明确的理论意义和工程价值。

\subsection{研究意义}
围绕交通资产清查中“资产发现不完整”与“异常状态记录不充分”两类核心问题，本文提出面向清查任务的主动感知检测与知识引导属性识别增强方法，其研究意义主要体现在以下三个方面。

\textbf{（1）理论意义：构建面向交通资产清查的整体任务分析框架。}本文不再将交通资产检测与属性识别视为彼此孤立的两个视觉任务，而是将其统一到交通资产清查这一上位应用目标下，形成“目标发现—状态理解—结构化记录”的整体任务链条。在此基础上，本文分别从主动感知、多智能体协作和知识引导生成三个角度，对交通资产清查中的关键机制进行重新组织，为复杂场景下的资产智能清查提供了较为系统的任务建模视角 \cite{agentthink2025,functional_agent}。

\textbf{（2）方法意义：分别突破漏检补偿与长尾异常识别两类关键瓶颈。}针对检测阶段的小目标与疑难区域漏检问题，本文通过引入主动感知流程和 Agent 功能解耦机制，构建“全局初检—候选区域推荐—局部精检—结果融合”的多阶段检测框架；针对属性识别阶段异常样本稀缺的问题，本文通过融合行业标准、专家经验和 Prompt 约束，构建知识引导的长尾数据生成与增强方案，以提升结构化异常属性识别能力 \cite{gb5768_2022,domain_knowledge_gen}。这两部分方法分别对应交通资产清查中的“查全资产”和“看懂状态”两项关键能力。

\textbf{（3）应用意义：提升交通资产清查与数字化养护的工程可用性。}面向交通资产清查的复合感知方法，能够在减少人工依赖的同时提升资产发现完整性和异常状态记录质量，从而为台账更新、风险研判和养护决策提供更可靠的数据基础。与仅关注单一检测精度的研究相比，本文更强调感知结果在交通资产管理流程中的可解释性、结构化和可落地性，这对于推动交通基础设施管理从经验驱动向数据驱动转型具有现实意义 \cite{omnisight_ai,roadvision_ai,vru_sensing2024}。

综上，本文研究并非单纯追求某一视觉模型的局部性能提升，而是立足交通资产清查这一实际应用需求，探索主动感知检测与知识引导属性识别增强相结合的方法体系，以提升交通资产清查的自动化、完整性和结构化水平。

\section{国内外研究现状}

为明确本文研究的学术定位与工程切入点，本节围绕交通资产清查任务相关的四个方向进行梳理：交通基础设施资产清查与检测算法研究、视觉语言模型在交通场景中的应用、智能体系统架构，以及生成式 AI 与长尾数据增强。通过分析现有研究进展与不足，可以更清楚地说明本文为何需要同时从主动感知检测与知识引导属性识别增强两方面展开研究。

\subsection{交通基础设施资产清查与检测算法研究}
在传统深度学习框架下，道路资产清查已经有较为系统的探索。Katsameniset al. 构建了DORIE（Dataset of Road Infrastructure Elements），面向护栏、路缘石、隔离设施等多类基础设施元素，系统评估了YOLO 等主流检测器在实际路网环境中的表现，认为小目标、遮挡和远距视角是当前模型的主要难点之一\cite{katsamenis2025dorie}。Kim et al. 提出的KRID（Korean Road Infrastructure Dataset）则覆盖高速公路、国道及城市道路34 类道路设施，基于YOLOv5、Faster R-CNN 等建立了识别基线，为后续基于深度学习的道路资产识别和维护决策提供了大规模数据支撑\cite{kim2025krid}。美国交通部门的研究也开始关注低成本采集与资产识别结合，例如Kuang et al. 开发手机端人工智能识别包，针对标线、路侧垃圾、交通标志和钢护栏等资产构建约5000 张样本的数据集，利用YOLO 系列模型实现移动采集条件下的自动识别和状态评估\cite{kuang2024mobile}。国内在路面病害与道路资产检测方面同样有大量工作。李宇宏et al. 基于VGG-16 分类与SSD 目标检测构建了车载视角路面病害综合检测系统，实现对坑槽、松散、车辙、裂缝等6 类病害的识别与ArcGIS 显示\cite{韩豫2023基于深度学习和}。傅蒙恩et al. 在YOLOv5 基础上改进锚框聚类和骨干网络，在四类路面缺陷检测任务中在基本不增加计算量的前提下显著提升了mAP 与FPS\cite{傅蒙恩2025基于}。近期还有工作围绕YOLOv8 轻量化改进，针对RDD2022 等开源病害数据集，对参数量和计算量进行大幅压缩，以满足车载设备的实时应用需求\cite{胡凤阔2024一种基于改进}。这些研究体现出深度学习检测模型在路面病害、交通标志及路侧附属设施自动识别方面的有效性，但总体上仍依赖于固定类别集和大规模标注数据，对新类别、罕见设施和复杂场景的适应性有限。

在此背景下，开放词汇（开放集）目标检测方法开始引入资产识别场景。Grounding DINO 将DINO 检测器与大规模图文“接地”预训练结合，通过将文本编码结果与视觉特征在Transformer 中深度融合，实现“文本提示+ 边界框”的开放集检测，并在COCO、LVIS 等数据集上显著提升了零样本检测性能\cite{liu2024grounding}。Ilyas et al. 系统评估了Grounding DINO 等开放词汇检测器在“非常规街景”中的表现，发现其虽然能凭借文本提示识别不少未见类别，但在分布外少见物体和复杂路侧环境下仍存在漏检与误检问题\cite{ilyas2024potential}。Zhu et al. 进一步在道路场景中提出DetSeg，将Grounding DINO 作为基础检测模块，引入时序信息和不确定性约束，对路侧场景中的分布外目标（如罕见障碍物）进行检测和分割，验证了开放词汇检测器在复杂交通场景安全监测中的潜力\cite{zhu2025detseg}。

部分工作已经开始将开放词汇检测直接用于交通基础设施及相关场景。例
如，国内专利中提出了“开放集交通基础设施三维数字化大模型”构建方案，利用基于文本提示的开放词汇检测模型结合SAM 分割和三维点云投影，实现对街景中多类交通基础设施的识别与三维重建\cite{CN120125811A}。在轨道交通领域，中科院自动化所的“紫东·太初”多模态大模型平台结合开放词汇目标检测，对高铁接触网异物进行识别，展示了开放词汇检测在铁路基础设施安全监测中的应用前景\cite{NCSTI2023AIrail}。产业界也出现了面向交通行业的开放词汇检测应用，例如臻识科技提出的VTraffic大模型，将开放词汇检测能力与交通事件识别结合，为视频卡口与路口场景的多目标识别提供支持\cite{7its2024trafficAImodel}。

整体来看，传统深度学习方法和开放词汇检测已经在资产识别层面形成了
较为丰富的研究基础。前者在封闭类别集上具有成熟的工程实践和较高的精度；后者通过引入文本编码和图文预训练，使模型在一定程度上具备了识别未见类别、扩展资产字典的能力。但现有工作大多以单帧图像的小尺度检测基准为依托，研究重点仍然集中在检测精度或开放类别能力本身，对车载/路侧高分辨率全景图像场景下的高质量感知问题关注较少，也很少与后端资产清查、台账系统进行深度联动。

\subsection{视觉语言模型相关研究及在交通场景中的应用}
在开放词汇检测的基础上，大规模视觉语言模型（VLM）的发展显著扩展
了多模态模型的能力边界。CLIP 通过在约4 亿对图文样本上进行对比学习，将图像编码器与文本编码器映射到统一特征空间，实现了无需任务特定标注即可进行零样本图像分类和检索的能力，在ImageNet 等基准上验证了其强大的开放集识别性能\cite{deng2009imagenet,radford2021learning}。随后，BLIP/BLIP-2 等工作采用“冻结视觉编码器+ 冻结大语言模型+ 轻量级跨模态模块”的结构，通过较小代价实现多模态问答与图文生成，并在多种视觉语言任务上取得优异结果\cite{li2022blip,li2023blip,ding2025urban}。在此基础上，LLaVA 等多模态大模型通过“视觉指令微调”构造大规模多模态指令数据，将视觉编码器与指令微调后的大语言模型连接，使模型具备了多轮对话、复杂视觉推理和细粒度描述能力，成为目前通用视觉助手的重要代表之一\cite{li2024llava,wang2024qwen2}。

在与交通基础设施密切相关的领域，VLM 的应用研究近两年发展较快。遥感与国土监测方向上，Efunogbon et al. 基于增强版RSCLIP 提出了面向遥感场景的大型视觉语言模型框架，讨论了模型在跨域适应、复杂地物属性表征方面的优势与难点\cite{efunogbon2025large}。Gao et al. 利用VLM 进行环境风险预测，将高分辨率遥感影像与文本说明结合，以提升复杂空间场景下的分类与解释能力，同时指出高分辨率输入与多尺度结构带来的计算开销问题\cite{gao2025vision}。

在道路与交通场景中，越来越多文献开始围绕“VLM + 路网/城市图像”展开。Ding et al. 提出URA-VLMs 框架，将InternVL 等基础多模态模型适配于城市道路异常监测，构建包含3000 余张城市路况图像的数据集，实现对路面积水、堆积物、拥堵等多类异常现象的识别与多维度安全评估，并引入RAG 机制对异常解释进行知识增强\cite{ding2025urban}。Zheng et al. 对基于VLM 的交通异常检测进行了系统综述，将现有方法归纳为基于提示（prompt-based）、端到端微调和适配器特征三大类，指出VLM 在复杂交通场景中对语义丰富的异常事件具有更强鲁棒性，但在时序建模和高分辨率细节感知方面仍有限\cite{zhengvideo}。Ren et al. 进一步提出CoT-VLM4Tar 框架，引入链式思维提示，引导VLM 在模拟交通环境中对异常事件进行多步推理与处置方案生成，展示了多模态链式思维在交通安全决策中的应用潜力\cite{ren2025cot}。

在道路损伤与基础设施方面，RoadBench 工作提出了首个面向道路病害理解的多模态基准，将路面裂缝、坑槽等损伤图像与详细文本描述配对，用于训练和评估视觉语言模型在“病害检测+ 语义解释”一体化任务中的表现\cite{xiao2025roadbench}。城市计算和自动驾驶社区也出现了利用CLIP/BLIP 等模型对驾驶场景进行开放词汇识别和场景检索的工作，例如通过自然语言描述查询特定交通场景、利用图–文相似度进行复杂场景检索等\cite{jia2024bev}，相关综述指出VLM 有望在自动驾驶环境认知与交通行为分析中发挥基础模型作用\cite{zhou2024vision}。

总体而言，这些研究表明VLM 在交通与道路基础设施相关任务中主要沿两个方向展开：一是围绕异常与风险识别，利用VLM 的开放词汇和推理能力统一表征多类异常现象；二是围绕病害与损伤解释，通过图文对或多模态指令数据，让模型同时输出检测结果及自然语言说明。不过，从现有文献来看，大多数工作仍采用中等分辨率图像或局部裁剪作为输入，对车载/路侧采集的大范围高分辨率全景影像，很少有专门针对“如何在可接受计算代价下保持对小目标和长尾资产感知能力”的系统研究，这一问题更多出现在近年的主动感知与高分辨率多模态框架中。

\subsection{智能体（Agent）系统架构}
随着视觉语言模型和大语言模型能力的快速提升，研究者逐渐从“单一模型完成检测任务”的范式，转向“多模块协同执行”的智能体（Agent）系统架构，使系统能够根据任务目标主动调用检测、分割、检索及知识库等工具，实现复杂场景下的闭环感知与决策。

在基础设施巡检领域，已有研究开始探索多智能体协作的巡检框架。例如，Yu et al. 针对隧道结构缺陷巡检提出了基于多模态大模型的双 Agent 系统，将缺陷识别 Agent 与报告生成 Agent 解耦：前者负责从巡检图像中检测裂缝、剥落等结构损伤，后者结合行业规范自动生成维护建议与巡检报告，实现了从视觉检测到运维决策的自动化闭环 \cite{yu2026tunneldualagent}。该研究表明，通过引入大模型作为认知中枢，可以显著提升基础设施巡检系统的自动化程度和信息整合能力。

在城市交通管理领域，智能体架构也被用于多源数据的融合与分析。Kalyuzhnaya et al. 提出面向智慧城市基础设施管理的 LLM Agent 系统，通过统一调度图像分析模块、数据库查询模块和规则推理模块，实现对城市事件和基础设施状态的自动识别与决策支持 \cite{kalyuzhnaya2025smartcity}。该系统能够根据任务目标动态选择所需工具链，并整合多源信息生成结构化结果，为城市资产管理提供新的系统实现路径。

在视觉感知系统工程方面，一些工业界方案已经采用“认知模型+视觉工具”的分层结构。例如 NVIDIA 提出的视觉分析 Agent 管线中，首先通过开放词汇检测和分割模块获取候选目标，再利用视觉语言模型对候选区域进行语义解释与事件总结，从而将原始视频转化为可查询的结构化信息 \cite{nvidia2025videoagents}。类似思路也被应用于交通基础设施检测，通过开放词汇检测器识别潜在目标，再结合多模态模型进行语义验证与分类，提高复杂场景下的识别可靠性。

与此同时，开放词汇检测研究也逐渐体现出与智能体架构结合的趋势。Zhu et al. 在道路场景中提出的 DetSeg 方法，将开放词汇检测作为候选目标生成模块，并结合不确定性约束进行后续筛选，有效降低了分布外目标检测中的误报率 \cite{zhu2025detseg}。这一类方法本质上已经形成了“候选生成—验证筛选”的多阶段处理流程，为进一步引入具备推理能力的认知 Agent 提供了基础。

总体来看，智能体架构正在成为交通基础设施巡检系统的重要发展方向。然而，现有研究多集中于缺陷报告生成或多源数据融合，针对车载或路侧高分辨率巡检图像中小目标资产的主动发现与漏检补偿机制仍缺乏系统研究，也尚未形成面向交通资产清查任务的完整多智能体检测框架。

\subsection{生成式 AI 与长尾数据增强}
交通基础设施资产清查中的一个关键难点在于受损状态和罕见资产的样本极度稀缺。相比完好设施，大面积锈蚀、断裂、倾斜或严重磨损等损坏状态在真实路网中出现频率较低，采集与标注成本高，导致训练数据呈现明显的长尾分布。这种数据不平衡问题会显著降低深度学习模型在受损资产识别中的泛化能力。

针对道路病害检测任务，研究者较早开始尝试通过合成数据缓解样本不足问题。Konushin et al. 通过在真实道路场景中嵌入合成交通标志图像，构建扩展数据集以提升稀缺类别识别能力，实验表明该方法能够显著提升交通标志检测模型的召回率 \cite{konushin2021syntheticsigns}。在路面裂缝检测领域，Cano et al. 利用基于扩散模型的语义条件生成方法，根据裂缝掩码生成高保真裂缝图像，从而扩展训练数据并提升模型在复杂裂缝形态下的识别能力 \cite{cano2024semanticdiffusion}。Hu et al. 提出的 RoadDiffBox 方法进一步将扩散模型生成与半监督学习结合，实现道路病害目标的受控生成和自动标注，有效降低了对人工标注数据的依赖 \cite{hu2025roaddiffbox}。

随着生成式模型能力的提升，研究开始关注交通设施受损状态的可控生成。Merolla et al. 利用 Stable Diffusion 等生成式模型合成交通标志损坏和路面病害图像，以缓解真实损坏样本不足的问题，并验证了生成数据在提升检测模型性能方面的有效性 \cite{merolla2024signsurfacedamage}。相关研究表明，通过合理的提示词设计和局部图像编辑，可以生成符合真实场景特征的损坏资产图像，为长尾类别提供有效补充。

在交通基础设施数字化管理领域，生成式 AI 也逐渐与资产识别系统结合。例如，在开放集交通基础设施三维数字化方案中，通过开放词汇检测识别资产后，结合语义模型对资产类别和状态进行扩展描述，为资产数字化管理提供更丰富的语义信息 \cite{CN120125811A}。产业界同样开始探索生成式 AI 在交通行业中的应用，将大模型能力用于交通场景理解与事件识别，为交通资产管理提供辅助支持 \cite{7its2024trafficAImodel}。

总体而言，生成式 AI 为解决交通资产清查中的长尾问题提供了新的技术路径。通过合成罕见类别和受损状态样本，可以显著缓解数据不平衡问题并提升模型泛化能力。然而，现有研究主要集中于道路病害或交通标志检测任务，对交通资产清查中多类型设施受损状态的系统化生成策略研究仍较少，尤其缺乏结合行业标准知识进行可控生成的数据增强方法。


\section{研究内容与目标}
围绕交通资产清查任务中“资产发现不完整”和“异常状态记录不充分”两类关键问题，本文从主动感知检测与知识引导属性识别增强两个层面展开研究，目标是构建一套面向交通资产清查的检测与异常属性识别方法体系，使系统既能够更完整地发现交通资产目标，又能够更准确地输出资产异常属性及结构化状态信息。

本文的主要研究内容包括以下三个方面：
\begin{enumerate}
    \item \textbf{面向交通资产清查的主动感知检测方法研究。}针对高分辨率道路场景中小目标、边缘目标和密集目标易漏检的问题，本文提出基于主动感知的多智能体交通资产检测框架，通过 Detection Agent 与 Crop Agent 的协同，将静态全图检测扩展为“全局初检—候选区域推荐—局部精检—结果融合”的三阶段流程，以提升资产发现完整性与漏检补偿能力。
    \item \textbf{面向交通资产属性识别增强的知识引导长尾数据生成方法研究。}针对交通资产异常状态样本稀缺、长尾分布显著的问题，本文提出知识引导的长尾数据生成与增强方法，将交通行业规范、专家经验与 Prompt 约束转化为可控生成条件，围绕关键异常缺口定向补齐长尾样本，并在统一多模态指令模型基座上验证其对结构化属性识别的促进作用。
    \item \textbf{面向清查应用的整体方法有效性验证。}本文分别从检测性能、异常属性识别性能和工程应用意义三个层面，对所提出方法进行系统实验分析，重点验证主动感知机制对资产查全能力的提升，以及知识引导生成对异常状态识别与结构化记录质量的改善效果。
\end{enumerate}

概括而言，本文的创新性主要体现在以下三个方面：其一，将交通资产检测组织为“全局初检—候选区域推荐—局部精检—结果融合”的主动感知流程，以提升资产查全能力；其二，将知识使用方式由被动辅助转变为面向长尾缺口的显式组织与条件约束，构建服务于属性识别增强的长尾数据生成方法；其三，从训练分布改善而非单纯模型放大的角度，系统验证知识引导长尾数据生成对结构化属性识别的促进作用。

基于上述研究内容，本文的总体目标可以概括为：构建一种面向交通资产清查的主动感知检测与知识引导属性识别增强方法，使交通资产清查过程在资产发现、异常识别和结构化记录三个关键环节上获得更高的自动化水平与更强的工程可用性。

\section{论文组织架构}

\textbf{第一章}为绪论，主要介绍论文的研究背景与研究意义，梳理交通资产清查、视觉语言模型、智能体系统以及生成式 AI 与长尾数据增强等相关研究现状，并进一步明确本文的研究内容、总体目标与论文整体结构。

\textbf{第二章}为相关理论方法与数据集介绍，围绕后续研究所需的理论基础展开说明，重点介绍目标检测任务定义与评价指标、常用目标检测方法、数据生成与长尾样本增强理论、多模态大模型识别与监督微调技术，并对 WHU-Infra3D 数据集及实验数据组织方式进行说明，为后续方法设计与实验分析提供理论与数据支撑。

\textbf{第三章}介绍基于主动感知的多智能体交通资产检测方法与实验。首先阐述基于异构 Agent 协作的三阶段检测框架、Detection Agent 与 Crop Agent 的功能分工以及 Crop Agent 的训练数据构建方法；随后介绍第三阶段局部精检与结果融合策略，并结合实验验证与结果分析，从整体性能、小目标与密集目标检测表现、消融实验等方面评估所提方法的有效性。

\textbf{第四章}介绍面向交通资产属性识别的长尾数据构建与增强验证。首先说明领域知识引导长尾数据生成的总体技术路线、属性识别增强的问题定义与研究动机，以及长尾异常样本构建与生成模型选择策略；随后在长尾数据增强的交通资产属性识别任务中开展实验验证与结果分析，重点评估知识引导长尾补样对整体属性识别、关键异常字段识别和结构化输出质量的提升作用。

\textbf{第五章}为总结与展望，对全文研究工作进行系统总结，归纳本文的主要创新点，并结合当前研究的局限性对后续可能的改进方向与进一步研究内容进行展望。
